\chapter{Contragredient Algebras and Virasoro}

All vector spaces in this chapter are over \(\mathbb C\).
The contragredient construction starts from a matrix and retains the
Cartan subalgebra while discarding the largest ideal disjoint from it.
This gives the structural meaning of the simple-root relations beyond
finite type. The second half constructs the Witt and Virasoro algebras
and computes their first highest-weight singular vectors.
For the loop bracket, the degree derivation and the concrete untwisted
affine highest-weight results we use \emph{Topics in Representation
Theory}, LO09--LO14 \cite[PDF pp.~26--28]{mg-reptopics}.

\section{Cartan data and the contragredient quotient}

\begin{definition}[Generalized Cartan matrices and realizations]
\label{res:KM01}
A generalized Cartan matrix \(A=(a_{ij})_{1\le i,j\le n}\) has
integer entries, \(a_{ii}=2\), \(a_{ij}\le0\) for \(i\ne j\), and
\(a_{ij}=0\) if and only if \(a_{ji}=0\). It is symmetrizable
if there are positive real numbers \(d_i\) such that
\(d_i a_{ij}=d_j a_{ji}\). A realization is a vector space
\(\mathfrak h\) with independent elements \(h_i\) and independent
linear forms \(\alpha_i\) such that
\(\alpha_j(h_i)=a_{ij}\). It is minimal when
\(\dim\mathfrak h=2n-\operatorname{rank}A\).
\end{definition}

\begin{lemma}[Existence and uniqueness of a minimal realization]
\label{res:KM02}
Every generalized Cartan matrix has a minimal realization. Any
two such realizations are related by an isomorphism carrying each
\(h_i\) to its counterpart and pulling back each counterpart of
\(\alpha_i\) to \(\alpha_i\). In a minimal realization,
\(\bigcap_i\ker\alpha_i\) has dimension \(n-\operatorname{rank}A\).
\end{lemma}

\begin{proof}
Let \(\mathfrak h_0\) have basis \(h_i\), and define
\(T_0:\mathfrak h_0\to\mathbb C^n\) by
\(T_0(h_i)=(a_{i1},\ldots,a_{in})\).
Choose a complement \(E\) of its image in \(\mathbb C^n\),
set \(\mathfrak h=\mathfrak h_0\oplus E\), and let
\(T(h,e)=T_0(h)+e\). The coordinate forms of \(T\) are independent
because \(T\) is onto. In any realization, the restriction of
\(T=(\alpha_i)_i\) to \(\mathfrak h_0\) has rank
\(\operatorname{rank}A\); surjectivity forces
\(\dim(\mathfrak h/\mathfrak h_0)\ge n-\operatorname{rank}A\).
For a minimal realization, choose a linear section of \(T\) over the
same complementary space \(E\). Its image is disjoint from
\(\mathfrak h_0\) and complements it by dimension. Thus the above
model identifies the realizations. The dimension of the common
kernel is rank-nullity for \(T\).
\end{proof}

\begin{proposition}[The universal algebra before the quotient]
\label{res:KM03}
Let \(\widetilde{\mathfrak g}\) be generated by the vector space
\(\mathfrak h\) and symbols \(e_i,f_i\), with relations
\[
 \begin{gathered}
 [h,h']=0,\qquad [h,e_i]=\alpha_i(h)e_i,\qquad
 [h,f_i]=-\alpha_i(h)f_i,\\
 [e_i,f_j]=\delta_{ij}h_i .
 \end{gathered}
\]
The displayed copy of \(\mathfrak h\) embeds,
and there is a vector space decomposition
\[
 \widetilde{\mathfrak g}
   =\widetilde{\mathfrak n}^-\oplus\mathfrak h
      \oplus\widetilde{\mathfrak n}^+ ,
\]
where the Lie subalgebras generated by the \(f_i\) and \(e_i\)
respectively are free on those generators. With
\(Q=\bigoplus_i\mathbb Z\alpha_i\), the grading has
\(\deg h=0\), \(\deg e_i=\alpha_i\), and
\(\deg f_i=-\alpha_i\). Each homogeneous space is the
corresponding weight space for the adjoint action of \(\mathfrak h\).
\end{proposition}

\begin{proof}
We use A05 and A07 of \emph{Lie Algebras and Finite-Dimensional
Representations} for PBW and the embedding of a free Lie algebra
in the corresponding tensor algebra \cite[PDF pp.~8--9]{mg-lie}.
Take the tensor space on symbols \(F_i\) and its vacuum \(v_\lambda\)
for any \(\lambda\in\mathfrak h^*\). On a word \(w\) let
\(\gamma(w)\) be the sum of the \(\alpha_j\) labeling its letters.
Let \(f_i\) prefix \(F_i\) and let \(h\) act diagonally on \(w\) by
\((\lambda-\gamma(w))(h)\). Define \(e_i\) to kill the vacuum and
recursively set
\[
 e_i(F_j w)=F_j(e_i w)
               +\delta_{ij}(\lambda-\gamma(w))(h_i)\,w.
\]
Every nonzero term of \(e_i w\) has \(\gamma\)-degree
\(\gamma(w)-\alpha_i\), by induction on the length. The weight
shift gives the required \(h,e_i\) bracket, and the recursion
gives exactly the \(e_i,f_j\) bracket. Thus these operators realize
the relations. If an \(h\) acts as zero for all \(\lambda\), its value
on each vacuum is zero, so \(h=0\). The action of words in the
\(f_i\) on the vacuum is free associative prefixing. Restricting
that action to the free Lie subalgebra proves freeness of
\(\widetilde{\mathfrak n}^-\). Interchanging \(e_i,f_i\) and sending
\(h\) to \(-h\) preserves the relations, hence gives freeness of
\(\widetilde{\mathfrak n}^+\) as well.

By induction on bracket length the relation for \([e_i,f_j]\)
gives \([e_i,\widetilde{\mathfrak n}^-]\subseteq
\mathfrak h+\widetilde{\mathfrak n}^-\): expand the bracket with
a Lie word into two terms by Jacobi. The induction terms
bracket elements of \(\mathfrak h\) or \(\widetilde{\mathfrak n}^-\)
with smaller negative words. The interchange gives the opposite
containment for \([f_i,\widetilde{\mathfrak n}^+]\).
The sum of the three subspaces contains the generators and is stable
under their adjoints. Iterated adjoints of the generators applied to
generators span the presented Lie algebra by Jacobi, so the sum
is all of it. The relations are homogeneous for \(Q\); the three
subspaces have only strictly negative, zero and strictly positive
degrees respectively. This gives directness. On a bracket word
the adjoint \(h\)-weight is its degree, by the derivation rule.
Independence of the \(\alpha_i\) identifies distinct degrees
with distinct weights.
\end{proof}

\begin{theorem}[The Cartan-disjoint ideal and the center]
\label{res:KM04}
There is a largest ideal \(\mathfrak r\) of
\(\widetilde{\mathfrak g}\) with \(\mathfrak r\cap\mathfrak h=0\).
The contragredient Lie algebra
\(\mathfrak g(A)=\widetilde{\mathfrak g}/\mathfrak r\) inherits
the triangular decomposition and the \(Q\)-grading.
Its center is \(\bigcap_i\ker\alpha_i\subset\mathfrak h\).
For a minimal realization its dimension is
\(n-\operatorname{rank}A\). Any ideal of \(\mathfrak g(A)\)
disjoint from \(\mathfrak h\) is zero.
\end{theorem}

\begin{proof}
Each ideal is a direct sum of its intersections with the weight
spaces. Indeed an element has only finitely many weight components.
Choose \(h\) separating their values, which is possible by avoiding
finitely many proper hyperplanes, and extract those components
as polynomial projections in \(\operatorname{ad}h\). Stability
under the adjoint keeps them in the ideal. Thus an ideal disjoint
from \(\mathfrak h\) has zero degree-zero component by
Proposition~\ref{res:KM03}. The sum of all these ideals still has
zero degree-zero component, and is the required largest ideal.
In the quotient, the triangular summands survive as the corresponding
degree ranges, including all of \(\mathfrak h\). A central vector is
of degree zero because it commutes with \(\mathfrak h\).
The \(e_i,f_i\) in the quotient are nonzero since
\([e_i,f_i]=h_i\ne0\); hence a vector of \(\mathfrak h\) commutes
with them precisely when all the \(\alpha_i\) vanish on it.
Lemma~\ref{res:KM02} gives the dimension. Finally the preimage of a
Cartan-disjoint ideal of the quotient is itself Cartan-disjoint and
contains \(\mathfrak r\), so it equals \(\mathfrak r\).
\end{proof}

\section{Simple reflections and highest weights}

\begin{theorem}[Serre relations in the quotient]
\label{res:KM05}
For \(i\ne j\), the simple generators of \(\mathfrak g(A)\) satisfy
\[
 (\operatorname{ad}e_i)^{1-a_{ij}} e_j=0,\qquad
 (\operatorname{ad}f_i)^{1-a_{ij}} f_j=0 .
\]
Each \(\operatorname{ad}e_i\) and \(\operatorname{ad}f_i\) is locally
nilpotent on \(\mathfrak g(A)\).
\end{theorem}

\begin{proof}
In \(\widetilde{\mathfrak g}\) write
\(v_s=(\operatorname{ad}e_i)^s e_j\) and \(m=1-a_{ij}\).
One has \([f_i,v_0]=0\) and \(v_s\) has \(h_i\)-weight
\(a_{ij}+2s\). Thus induction gives
\[
 [f_i,v_s]=-s(a_{ij}+s-1)v_{s-1}\qquad(s\ge1),
\]
by expanding \([f_i,[e_i,v_{s-1}]]\) using Jacobi and
\([f_i,e_i]=-h_i\). At \(s=m\) this is zero.
For \(k\notin\{i,j\}\) the adjoint of \(f_k\) commutes with
\(\operatorname{ad}e_i\) and kills \(e_j\), so it kills \(v_m\).
The \(f_j\)-bracket is \(-(\operatorname{ad}e_i)^m h_j\).
It vanishes when \(m\ge2\), since the first bracket is already a
multiple of \(e_i\). When \(m=1\), \(a_{ij}=0\) gives
\(a_{ji}=0\), so it still vanishes.

Take the span of \(v_m\) and all successive adjoints of the
\(e_k\) on it. It is stable under \(\mathfrak h\) by homogeneity,
and under the \(e_k\) by construction. It is stable under each
\(\operatorname{ad}f_k\), by commuting this operator past the
first \(\operatorname{ad}e_\ell\) and using
\([f_k,e_\ell]=-\delta_{k\ell}h_k\); induction uses that \(f_k\)
kills \(v_m\). This span is therefore an ideal, with strictly
positive degrees, and lies in the ideal of
Theorem~\ref{res:KM04}. The interchange \(e_i\leftrightarrow f_i\),
\(h\mapsto-h\), preserves \(\mathfrak r\) by its maximality, and
gives the negative relations. On the simple generators, local
nilpotence of an adjoint \(\operatorname{ad}e_i\) follows from
these relations, from \([e_i,f_j]=\delta_{ij}h_i\), and from
\([e_i,h]=-\alpha_i(h)e_i\). The derivation rule makes it
locally nilpotent on every bracket word, because a sufficiently
high iterate then hits some generator beyond its bound.
The negative case follows by interchange.
\end{proof}

\begin{proposition}[Reflection automorphisms and real roots]
\label{res:KM06}
For each \(i\) there is an automorphism
\[
 r_i=\exp(\operatorname{ad}e_i)
      \exp(-\operatorname{ad}f_i)\exp(\operatorname{ad}e_i)
\]
whose restriction to \(\mathfrak h\) sends
\(h\) to \(h-\alpha_i(h)h_i\).
It takes the root space of \(\beta\) isomorphically to that
of \(s_i\beta=\beta-\beta(h_i)\alpha_i\). All roots in the
orbits of the simple roots under the \(s_i\) have multiplicity one.
If \(A\) is symmetrizable, the form on the real span of the
root lattice with
\((\alpha_i,\alpha_j)=d_i a_{ij}\) is invariant under each
\(s_i\). In this case a root in the orbit of \(\alpha_i\)
has squared length \(2d_i\), even when the form is degenerate.
\end{proposition}

\begin{proof}
All exponentials are finite sums on each vector by
Theorem~\ref{res:KM05}. The iterated derivation rule for brackets
shows \(\exp(D)[x,y]=[\exp(D)x,\exp(D)y]\); the inverse is
\(\exp(-D)\), by the binomial identity. In the three factors above,
the first applied to \(h\) gives \(h-\alpha_i(h)e_i\),
the second sends \(h\) to \(h-\alpha_i(h)f_i\) and \(e_i\) to
\(e_i+h_i-f_i\). Applying the last factor gives \(h-\alpha_i(h)h_i\),
using \([h_i,e_i]=2e_i\). This map on \(\mathfrak h\) is an
involution, so the new weight of \(r_i x\), for \(x\) of
weight \(\beta\), is \(\beta\) composed with that map. The
\(\alpha_i\) root space is the line \(\mathbb C e_i\): a bracket
word of that degree uses just \(e_i\), and \(e_i\ne0\).
Its orbit thus has the indicated multiplicities. The formula
\(\beta(h_i)= (\alpha_i,\beta)/d_i\) expresses \(s_i\) as
reflection in \(\alpha_i\) for the symmetric form. Direct
expansion preserves the form since
\((\alpha_i,\alpha_i)=2d_i\); no inverse of the form is used.
\end{proof}

\begin{example}[An infinite rank-two root orbit]
\label{res:KM07}
For \(a,b>0\) integers with \(ab\ge4\), let
\(A=\left(\begin{smallmatrix}2&-a\\-b&2\end{smallmatrix}\right)\).
The algebra \(\mathfrak g(A)\) is infinite-dimensional.
\end{example}

\begin{proof}
In simple-root coordinates the matrix of \(s_1s_2\) is
\(\left(\begin{smallmatrix}ab-1&-a\\ b&-1\end{smallmatrix}\right)\),
of determinant \(1\) and trace \(ab-2\). If \(ab>4\), its
two eigenvalues are positive reciprocal real numbers, different
from \(1\). Thus no nonzero vector is fixed by any positive
power of this matrix. If \(ab=4\), the matrix is \(I+N\) with
\(N^2=0\) by its characteristic polynomial, and
\(N\alpha_1=2\alpha_1+b\alpha_2\ne0\).
Then \((I+N)^j\alpha_1=\alpha_1+jN\alpha_1\) are all distinct.
In either case Proposition~\ref{res:KM06} gives infinitely many
nonzero, distinct root spaces in the orbit of \(\alpha_1\).
\end{proof}

\begin{theorem}[Highest-weight modules and integrability]
\label{res:KM08}
For \(\lambda\in\mathfrak h^*\), the Verma module induced from the
one-dimensional \(\mathfrak h\oplus\mathfrak n^+\)-module of
weight \(\lambda\) has finite-dimensional weight spaces,
one-dimensional top, and a unique simple quotient \(L(\lambda)\).
On this simple quotient all \(e_i,f_i\) are locally nilpotent
if and only if \(\lambda(h_i)\in\mathbb Z_{\ge0}\) for all \(i\).
In that case the multiplicities of its weight spaces are invariant
under the reflections \(s_i\) of Proposition~\ref{res:KM06}.
\end{theorem}

\begin{proof}
PBW, Result A05 of \cite[PDF pp.~8--9]{mg-lie}, for an order
with the negative-root basis first gives the
Verma module as \(U(\mathfrak n^-)\) applied freely to the
highest vector. Every nonzero degree in the negative algebra
is a strictly negative combination of the simple degrees.
For a given weight difference \(\gamma=\sum_i m_i\alpha_i\)
with \(m_i\ge0\), its negative enveloping component is spanned
by words in the \(f_i\) with exactly \(m_i\) occurrences of
each \(i\), of which there are finitely many.
The degree of the vacuum is zero, and every proper submodule
misses the vacuum. Such submodules are sums of weight components,
by the polynomial extraction used in Theorem~\ref{res:KM04}.
Their sum is therefore still proper and is the unique maximal
submodule.

On a highest vector \(v\) the rank-one relation gives
\[
 e_i f_i^r v
   =r(\lambda(h_i)-r+1)f_i^{r-1}v
\]
by induction, as in the affine criterion LO12
\cite[PDF p.~27]{mg-reptopics}. If \(r>0\) is minimal with
\(f_i^r v=0\), its coefficient must vanish, proving necessity.
Conversely, for \(m=\lambda(h_i)\in\mathbb Z_{\ge0}\),
the vector \(f_i^{m+1}v\) in the Verma module is killed by \(e_i\)
by the formula, and by every other \(e_j\) since \([e_j,f_i]=0\).
It generates a proper submodule: PBW shows that the generated
weights are at most its own lower weight. It therefore vanishes
in the simple quotient. Each \(e_i\) already kills its highest
vector. For \(x=e_i\) or \(f_i\), Theorem~\ref{res:KM05} and the
derivation rule give local nilpotence of \(\operatorname{ad}x\)
on \(U(\mathfrak g(A))\). Now
\[
 x^N(uv)=\sum_{j=0}^N\binom Nj
       (\operatorname{ad}x)^j(u)x^{N-j}v
\]
proves local nilpotence on the entire module by taking \(N\)
large enough for both bounds. On an integrable module the
products \(\exp(e_i)\exp(-f_i)\exp(e_i)\) act invertibly.
The same binomial formula shows that conjugation by
\(\exp(x)\) on the action of an algebra element is
the action of \(\exp(\operatorname{ad}x)\). In particular the
product intertwines the action of \(\mathfrak h\) with that
of its image under \(r_i\). It therefore takes the weight
space of \(\mu\) isomorphically to that of \(s_i\mu\).
\end{proof}

\section{Comparison with the untwisted affine model}

\begin{proposition}[The contragredient affine algebra is the loop model]
\label{res:KM09}
Let \(\mathfrak a\) be finite-dimensional complex simple and take
the untwisted affine loop algebra with its degree derivation
as in LO09--LO10. The simple forms
\(\alpha_0=\delta-\theta,\alpha_1,\ldots,\alpha_\ell\) and
coroots \(h_0=K-h_\theta,h_1,\ldots,h_\ell\) give a
minimal realization of a symmetrizable generalized Cartan
matrix. The associated \(\mathfrak g(A)\) is the full loop
model, with the center and the derivation kept distinct.
\end{proposition}

\begin{proof}
Use the construction and conventions of LO09--LO10
\cite[PDF p.~26]{mg-reptopics}. Here \(\theta\) is the highest
finite root. The simple forms are independent since
\(\alpha_0(d)=1\) and all finite simple forms vanish on \(d\).
The coroots are independent by the coefficient of \(K\).
The finite simple Cartan block is invertible. The full matrix
has rank \(\ell\): the nonzero element \(K\) in the span of the
coroots is annihilated by all simple forms, and the finite block
gives the lower bound. Thus the realization of dimension
\(\ell+2\) is minimal. In the entries involving \(0\) and a
finite simple index, the two off-diagonal entries are
\(-\alpha_i(h_\theta)\) and \(-\theta(h_i)\), both nonpositive
integers by the root strings and dominance of \(\theta\), and
both zero exactly when \((\theta,\alpha_i)=0\), using D04 and
E09 of \cite[PDF pp.~24--25, 34--35]{mg-lie}.
The \(0,0\) entry is \(2\). Extending the finite root form
by making \(\delta\) orthogonal to everything gives
\(a_{ij}=2(\alpha_i,\alpha_j)/(\alpha_i,\alpha_i)\).
Thus the positive half-squared lengths symmetrize the full matrix.

The simple triples and the affine Cartan subalgebra satisfy the
relations of Proposition~\ref{res:KM03}, and generate the loop
model by LO10. The resulting surjection is injective on
\(\mathfrak h\). To see that its kernel is the largest
Cartan-disjoint ideal, any nonzero ideal of the loop model
contains a nonzero weight vector by polynomial weight
extraction. A nonzero real root vector
\(x_\alpha\otimes t^m\) brackets with a suitable opposite
root vector in degree \(-m\) to give a nonzero Cartan
element: the finite component is a nonzero multiple of
\(h_\alpha\). An imaginary degree vector
\(h\otimes t^m\) with \(m\ne0\) brackets with some
\(h'\otimes t^{-m}\) to give
\(m B(h,h')K\ne0\), since the invariant form is nondegenerate
on the finite Cartan. A nonzero element of weight zero
already lies in the affine Cartan. The largest ideal
\(\mathfrak r\) of Theorem~\ref{res:KM04} contains the kernel,
and its image has zero intersection with the affine Cartan:
lift a purported Cartan element along the isomorphism on
\(\mathfrak h\) and subtract that lift to land in the
kernel. Thus the image of \(\mathfrak r\) is zero, proving
the comparison.
\end{proof}

\section{Witt and its universal central extension}

\begin{proposition}[The Witt algebra is simple and perfect]
\label{res:KM10}
Every derivation of \(\mathbb C[t,t^{-1}]\) is a finite linear
combination of the elements \(L_m=-t^{m+1}\partial_t\),
\(m\in\mathbb Z\). These form the Witt algebra \(W\) with bracket
\([L_m,L_n]=(m-n)L_{m+n}\). The Lie algebra \(W\) is simple
and equals its derived algebra.
\end{proposition}

\begin{proof}
A derivation is determined by its value at \(t\), because the
Leibniz rule gives its values on \(t^n\) for positive powers
and \(D(t^{-1})=-t^{-2}D(t)\). The commutator of the
displayed derivations is computed at \(t\), giving the bracket.
The eigenvalues \(-m\) of \(\operatorname{ad}L_0\) are distinct.
Every element has finite support, so any nonzero ideal contains
some \(L_m\) by polynomial weight extraction. If \(m\ne0\)
then \([L_m,L_{-m}]=2mL_0\) lies in the ideal. With \(L_0\)
in the ideal, all \(L_n\), \(n\ne0\), are obtained by brackets.
The same two bracket formulas show that \(W=[W,W]\).
\end{proof}

\begin{theorem}[The Virasoro cocycle and the universal extension]
\label{res:KM11}
For an arbitrary complex vector space \(Z\) with trivial \(W\)-action,
every central \(Z\)-valued Lie \(2\)-cocycle of \(W\), modulo
coboundaries, has a unique representative of the form
\[
 c(L_m,L_n)=\delta_{m,-n}\frac{m^3-m}{12}\,z,\qquad z\in Z.
\]
In particular \(H^2(W,\mathbb C)\) has dimension one.
The Lie algebra \(\operatorname{Vir}=W\oplus\mathbb C C\)
with bracket
\[
 [L_m,L_n]=(m-n)L_{m+n}
             +\delta_{m,-n}\frac{m^3-m}{12}C,\qquad [C,L_m]=0 ,
\]
is perfect, has center \(\mathbb C C\), and is the universal
central extension of \(W\).
\end{theorem}

\begin{proof}
For trivial coefficients a cocycle is an alternating bilinear map
satisfying
\[
 c([x,y],u)+c([y,u],x)+c([u,x],y)=0 .
\]
Replacing it by \(c-\phi\circ[-,-]\) corresponds to changing
the linear section of a central extension by \(\phi\).
For \(j\ne0\) choose
\(\phi(L_j)=-c(L_0,L_j)/j\), so that after the change
\(c(L_0,L_j)=0\) for every \(j\). The cocycle equation on
\((L_0,L_m,L_n)\) now gives
\((m+n)c(L_m,L_n)=0\). Set \(b_m=c(L_m,L_{-m})\).
A further change with \(\phi\) supported at \(L_0\) and
\(\phi(L_0)=b_1/2\) makes \(b_1=0\) and retains the first
normalization. Then \(b_{-m}=-b_m\). On
\((L_m,L_n,L_{-m-n})\) the cocycle equation is
\[
 (m-n)b_{m+n}-(m+2n)b_m+(2m+n)b_n=0 .
\]
With \(n=1\) and \(m\ge2\) this gives
\((m-1)b_{m+1}=(m+2)b_m\). The unique solution for \(m\ge2\)
is \(b_m=(m^3-m)b_2/6\); the normalized values for
\(m=0,\pm1\) and antisymmetry give all integers. Conversely
those values satisfy the three-term relation for every \(m,n\),
as follows by expanding the cubic polynomial, and the cocycle
equation is zero off total index zero. This proves existence
and the stated form with \(z=2b_2\). A coboundary preserving
the first normalization has \(\phi(L_j)=0\) for \(j\ne0\),
and preserving \(b_1=0\) forces \(\phi(L_0)=0\).
Thus the normalized representative is unique.

The displayed bracket on \(\operatorname{Vir}\) is therefore
Lie. Its center is exactly \(\mathbb C C\) by simplicity and
noncommutativity of \(W\). In \(\operatorname{Vir}\), all
\(L_m\) and \(C\) are brackets or linear combinations thereof:
use first the brackets of \(L_0\) with \(L_n\), then
\([L_1,L_{-1}]=2L_0\) and
\([L_2,L_{-2}]=4L_0+C/2\). For any central extension
\(E\to W\), choose a linear section and form the cocycle with
values in its central kernel \(Z\).
After normalization, the chosen lifts of \(L_m\) and the vector
\(z\in Z\) define a Lie homomorphism
\(\operatorname{Vir}\to E\) over \(W\) by sending \(C\) to \(z\).
Two such homomorphisms differ by a central-valued linear map
and so agree on every bracket. Perfectness implies their equality.
This is the asserted universal property.
\end{proof}

\begin{proposition}[Witt derivations lift to the affine central extension]
\label{res:KM12}
In the loop model from Proposition~\ref{res:KM09}, before adjoining
the degree derivation, write \(x_i=x\otimes t^i\).
The assignments
\[
 D_m(x_i)=-i x_{m+i},\qquad D_m(K)=0
\]
give an action of \(W\) by derivations on the central extension
of the loop algebra. They therefore give a semidirect Lie algebra
with \(\operatorname{Vir}\) acting through \(W\); in this action
\(D_0\) is the negative of the loop degree derivation.
\end{proposition}

\begin{proof}
For the loop part, \(D_m\) applies \(L_m\) to the Laurent-polynomial
factor, so it is a derivation and the commutators give the Witt
action of Proposition~\ref{res:KM10}. For \(x_i,y_j\), the central
terms on the right side of the derivation identity sum to
\[
 \bigl(-i(m+i)-ji\bigr)\delta_{m+i+j,0}B(x,y)K=0 .
\]
This uses exactly the old loop cocycle of LO09. The central
term on the left side also vanishes since \(D_m(K)=0\).
The commutator equalities extend to \(K\), where all operators
vanish. Composing with the quotient
\(\operatorname{Vir}\to W\) gives an action of \(\operatorname{Vir}\);
the semidirect bracket satisfies Jacobi by the derivation and
commutator identities. The degree formula is immediate from
\(D_0(x_i)=-i x_i\).
\end{proof}

\section{Virasoro highest weights}

\begin{theorem}[Virasoro Verma modules]
\label{res:KM13}
For \(c,h\in\mathbb C\), there is a universal cyclic
highest-weight module \(M(c,h)\) with vector \(v\) satisfying
\(Cv=cv\), \(L_0v=hv\), \(L_jv=0\) for \(j>0\).
It has a basis
\[
 L_{-i_1}\cdots L_{-i_s}v,\qquad
       1\le i_1\le\cdots\le i_s,\quad s\ge0 .
\]
The weight of this vector for \(L_0\) is \(h+i_1+\cdots+i_s\).
Consequently the graded dimension with the top in degree zero is
\(\prod_{i\ge1}(1-q^i)^{-1}\) as a formal power series,
and \(M(c,h)\) has a unique simple quotient.
\end{theorem}

\begin{proof}
Induce the one-dimensional module of
\(\mathbb C C\oplus\bigoplus_{j\ge0}\mathbb C L_j\)
with the specified action. The bracket relations of
Theorem~\ref{res:KM11} show that this one-dimensional action
is well-defined. Induction gives universality. PBW, Result A05 of
\cite[PDF pp.~8--9]{mg-lie}, for a total order with the negative
generators first gives the indicated basis.
The bracket \([L_0,L_{-i}]=iL_{-i}\) gives the weight.
Choosing the number of occurrences of each \(i\) gives the product
of geometric series, and each degree has only finitely many such
choices. A proper submodule misses the top since the top generates,
and it splits into weight spaces by polynomial extraction.
The sum of all proper submodules thus still misses the top
and gives the unique simple quotient.
\end{proof}

\begin{example}[Singular vectors in degrees one and two]
\label{res:KM14}
In \(M(c,h)\), a nonzero vector of degree \(1\) killed by all
\(L_j\) with \(j>0\) exists exactly when \(h=0\). In degree \(2\)
such a vector exists exactly when
\[
 18h-(4h+2)(4h+c/2)=0 .
\]
In that case the space of these vectors is the line spanned by
\[
 3 L_{-1}^{\,2}v-(4h+2)L_{-2}v .
\]
Both kinds of vectors generate proper submodules. In particular,
at \((c,h)=(1,1/4)\) the first such degree is \(2\), where the
vector is proportional to \(L_{-1}^{\,2}v-L_{-2}v\).
\end{example}

\begin{proof}
In degree \(1\) the basis is \(L_{-1}v\) and \(L_1L_{-1}v=2hv\).
Higher \(L_j\) kill this vector by their weight shifts.
In degree \(2\), take the ordered basis
\((L_{-2}v,L_{-1}^2v)\). The maps \(L_1,L_2\) to the
degrees \(1,0\), respectively, have rows
\[
 (3,4h+2), \qquad (4h+c/2,6h).
\]
For instance \([L_1,L_{-2}]=3L_{-1}\);
commuting \(L_1\) past both factors of \(L_{-1}^2\)
gives \(2(L_0L_{-1}+L_{-1}L_0)v=(4h+2)L_{-1}v\).
Similarly \([L_2,L_{-2}]=4L_0+C/2\) and commuting
\(L_2\) past both factors of \(L_{-1}^2\) gives \(6hv\).
The maps \(L_j\) for \(j\ge3\) again kill this space by degree.
The determinant of the two rows gives the criterion, and since
the first entry of the first row is \(3\), its nullspace is the
displayed line when the determinant is zero.
PBW moves all positive generators to the right of a word
applied to a singular vector, where they kill it, and moves
the Cartan factors to act as scalars. Thus the submodule generated
in either positive degree cannot contain the original top.
Substituting \(c=1,h=1/4\) gives the example.
\end{proof}

The contragredient construction and root symmetries are compared with
Kac, Chapter 1 and Sections 3.1--3.7 \cite{kac}.
The untwisted affine convention agrees with Hernandez,
Sections 3--4 \cite{hernandez-affine}. For Virasoro and highest
weights, see Kac, Sections 7.3 and 9.14 \cite{kac}.
The converse assertion that the Serre ideal already equals the
Cartan-disjoint ideal in symmetrizable type, the Weyl--Kac character
formula, and the higher-degree classification of Virasoro singular
vectors are further topics.
